% 第2章 2.1.1：算子范数作为单位球像的最大伸长。
% A=[[3/2,1/2],[1/2,3/2]]，特征值 2（方向 (1,1)/sqrt2）与 1（方向 (1,-1)/sqrt2）；
% 单位圆的像是半轴 2、1 的椭圆（旋转 45 度）。e_1=(1,0) 映为 (3/2,1/2)，长度 sqrt(5/2)≈1.58；
% v=(1,1)/sqrt2 映为 (sqrt2,sqrt2)，长度 2=‖A‖。比例：1 单位 = 9mm。
\begin{tikzpicture}[x=9mm,y=9mm,font=\figurefont\footnotesize,text=FigureInk,
  note/.style={font=\FigureNoteFont\fontsize{7.5}{9}\selectfont,text=FigureMuted},
  ax/.style={draw=FigureStroke,line width=.6pt,-{Stealth[length=1.6mm]}},
  vec/.style={line width=1pt,-{Stealth[length=2mm]}}]
\begin{scope}
  \draw[ax] (-1.6,0)--(1.7,0); \draw[ax] (0,-1.6)--(0,1.7);
  \fill[FigureBlueFill] (0,0) circle (1);
  \draw[FigureBlue,line width=1.1pt] (0,0) circle (1);
  \draw[vec,FigureGreen] (0,0)--(1,0) node[below right=-.5mm] {$\bm e_1$};
  \draw[vec,FigureYellow] (0,0)--(0.7071,0.7071) node[above right=-.6mm] {$\bm v$};
  \node[anchor=west] at (-1.9,2.0) {(a) 单位圆 $\|\bm x\|_2\leq1$};
\end{scope}
\draw[-{Stealth[length=2mm]},FigureStroke,line width=.8pt] (2.1,0.2) to[bend left=20] node[above] {$\bm A$} (3.3,0.2);
\begin{scope}[xshift=55mm]
  \draw[ax] (-2.3,0)--(2.4,0); \draw[ax] (0,-2.3)--(0,2.4);
  \fill[FigureBlueFill,rotate=45] (0,0) ellipse (2 and 1);
  \draw[FigureBlue,line width=1.1pt,rotate=45] (0,0) ellipse (2 and 1);
  \draw[densely dashed,FigureMuted,line width=.6pt] (0,0) circle (2);
  \draw[vec,FigureGreen] (0,0)--(1.5,0.5) node[right] {$\bm A\bm e_1$};
  \draw[vec,FigureYellow] (0,0)--(1.4142,1.4142) node[above right=-.6mm] {$\bm A\bm v$};
  \node[note,anchor=west] at (-2.3,-2.55) {虚线圆半径 $2=\|\bm A\|$};
  \node[anchor=west] at (-2.6,2.75) {(b) 像 $\bm A\{\|\bm x\|_2\leq1\}$};
\end{scope}
\end{tikzpicture}
